% 第3章习题解答
% 基于 chapters/ch03.tex；按原章题目出现顺序编排。
\chapter{第3章 实数系的基本定理习题解答}

\section{确界的概念和确界存在定理}
\subsection*{3.1.3 练习题}

\begin{xexercise}[3.1.3 练习题 1]
试证明确界的唯一性。
\end{xexercise}
\begin{xsolution}
先证上确界。设 $\alpha$ 与 $\beta$ 都是数集 $A$ 的上确界。因为 $\alpha$ 是最小上界，而 $\beta$ 是 $A$ 的一个上界，所以 $\alpha\leqslant\beta$；交换 $\alpha,\beta$ 的地位，又因 $\beta$ 是最小上界而 $\alpha$ 是上界，得 $\beta\leqslant\alpha$。故 $\alpha=\beta$。

再证下确界。若 $\gamma,\delta$ 都是 $A$ 的下确界，则 $\gamma$ 是最大下界而 $\delta$ 是一个下界，所以 $\gamma\geqslant\delta$；同时 $\delta$ 是最大下界而 $\gamma$ 是一个下界，所以 $\delta\geqslant\gamma$。因此 $\gamma=\delta$。故上下确界若存在都唯一。
\end{xsolution}

\begin{xexercise}[3.1.3 练习题 2]
设对每个 $x\in A$ 成立 $x<a$。问：在 $\sup A<a$ 和 $\sup A\leqslant a$ 中哪个是对的？
\end{xexercise}
\begin{xsolution}
必然成立的是
\[
  \sup A\leqslant a.
\]
因为 $x<a$ 对一切 $x\in A$ 成立，所以 $a$ 是 $A$ 的上界，最小上界不超过任一上界。

一般不能推出 $\sup A<a$。例如取 $A=(0,a)$（设 $a>0$），则对每个 $x\in A$ 都有 $x<a$，但 $\sup A=a$。
\end{xsolution}

\begin{xexercise}[3.1.3 练习题 3]
设数集 $A$ 以 $\beta$ 为上界，又有数列 $\{x_n\}\subset A$ 和 $\lim_{n\to\infty}x_n=\beta$。证明：$\beta=\sup A$。
\end{xexercise}
\begin{xsolution}
因为 $\beta$ 是 $A$ 的上界，只需证明任何小于 $\beta$ 的数都不是 $A$ 的上界。任取 $c<\beta$，令
\[
  \varepsilon=\frac{\beta-c}{2}>0.
\]
由 $x_n\to\beta$，存在 $N$，当 $n>N$ 时有
\[
  x_n>\beta-\varepsilon=\frac{\beta+c}{2}>c.
\]
而 $x_n\in A$，故 $c$ 不是 $A$ 的上界。于是 $\beta$ 是 $A$ 的最小上界，即 $\beta=\sup A$。
\end{xsolution}

\begin{xexercise}[3.1.3 练习题 4]
求原题所列七个数集的上确界和下确界。
\end{xexercise}
\begin{xsolution}
逐项计算如下。
\begin{enumerate}[label=(\arabic*)]
  \item 对 $A=\{x\in\QQ\mid x>0\}$，任意正实数都不是上界，故 $A$ 无上界；而 $0$ 是下界，且任意 $\alpha>0$ 都不是下界。因此
  \[
    \sup A=+\infty,\qquad \inf A=0.
  \]

  \item 当 $x\in(-\frac12,1)$ 时，$0\leqslant x^2<1$，且 $x=0$ 时取得 $0$；另一方面令 $x\to1^-$ 可使 $x^2$ 任意接近 $1$。故
  \[
    \inf\{x^2\mid x\in(-\tfrac12,1)\}=0,\qquad
    \sup\{x^2\mid x\in(-\tfrac12,1)\}=1.
  \]

  \item 记 $a_n=(1+1/n)^n$。由二项式展开
  \[
    a_n=\sum_{k=0}^{n}\frac1{k!}\prod_{j=0}^{k-1}\left(1-\frac{j}{n}\right).
  \]
  对每个固定的 $k\geqslant2$，乘积随 $n$ 增大而增大，且从 $n$ 到 $n+1$ 还增加一个正项，因此 $\{a_n\}$ 严格增加。又由已知基本极限 $a_n\to\ee$，故
  \[
    \inf\{a_n\}=a_1=2,
    \qquad
    \sup\{a_n\}=\ee.
  \]

  \item 令 $a_n=n\ee^{-n}>0$。有
  \[
    \frac{a_{n+1}}{a_n}=\left(1+\frac1n\right)\ee^{-1}<1,
  \]
  因为 $1+1/n\leqslant2<\ee$。故 $a_n$ 严格减少，而 $a_n\to0$。因此
  \[
    \sup\{n\ee^{-n}\}=\ee^{-1},\qquad
    \inf\{n\ee^{-n}\}=0.
  \]

  \item $\arctan x$ 在 $\RR$ 上取值于 $(-\pi/2,\pi/2)$，且当 $x\to\pm\infty$ 时分别趋于 $\pm\pi/2$。故
  \[
    \sup\{\arctan x\mid x\in\RR\}=\frac\pi2,
    \qquad
    \inf\{\arctan x\mid x\in\RR\}=-\frac\pi2.
  \]

  \item 令
  \[
    a_n=(-1)^n+\frac1n(-1)^{n+1}=(-1)^n\left(1-\frac1n\right).
  \]
  偶数项为 $1-1/n\to1$，奇数项为 $-1+1/n\to-1$，且所有项严格位于 $(-1,1)$ 内。因此
  \[
    \sup\{a_n\}=1,
    \qquad
    \inf\{a_n\}=-1.
  \]

  \item 对 $a_n=1+n\sin(n\pi/2)$，当 $n=4k+1$ 时 $a_n=1+n\to+\infty$；当 $n=4k+3$ 时 $a_n=1-n\to-\infty$。故
  \[
    \sup\{a_n\}=+\infty,
    \qquad
    \inf\{a_n\}=-\infty.
  \]
\end{enumerate}
\end{xsolution}

\begin{xexercise}[3.1.3 练习题 5]
证明原题中的两个确界不等式。
\end{xexercise}
\begin{xsolution}
设有关确界均有意义。
\begin{enumerate}[label=(\arabic*)]
  \item 对每个 $n$，有
  \[
    x_n\leqslant\sup\{x_n\},\qquad
    y_n\leqslant\sup\{y_n\},
  \]
  因而
  \[
    x_n+y_n\leqslant\sup\{x_n\}+\sup\{y_n\}.
  \]
  右端是数集 $\{x_n+y_n\}$ 的上界，于是
  \[
    \sup\{x_n+y_n\}\leqslant\sup\{x_n\}+\sup\{y_n\}.
  \]

  \item 对下确界，由定义对每个 $n$ 有
  \[
    x_n+y_n\geqslant\inf\{x_n\}+\inf\{y_n\},
  \]
  故右端是 $\{x_n+y_n\}$ 的下界，从而
  \[
    \inf\{x_n+y_n\}\geqslant\inf\{x_n\}+\inf\{y_n\}.
  \]
\end{enumerate}
\end{xsolution}

\begin{xexercise}[3.1.3 练习题 6]
设对任意 $x\in A$ 与任意 $y\in B$ 都有 $x<y$。证明：$\sup A\leqslant\inf B$。
\end{xexercise}
\begin{xsolution}
任取 $y_0\in B$，则 $y_0$ 是 $A$ 的上界，因此 $A$ 有上确界。任取 $x_0\in A$，则 $x_0$ 是 $B$ 的下界，因此 $B$ 有下确界。

对任意 $y\in B$，由于 $y$ 是 $A$ 的上界，有 $\sup A\leqslant y$。这说明 $\sup A$ 是 $B$ 的一个下界，故
\[
  \sup A\leqslant\inf B.
\]
\end{xsolution}

\begin{xexercise}[3.1.3 练习题 7]
设 $B=\{x+c\mid x\in A\}$。证明 $\sup B=\sup A+c$，$\inf B=\inf A+c$。
\end{xexercise}
\begin{xsolution}
先证上确界。记 $\alpha=\sup A$。对任意 $y\in B$，有 $y=x+c$，其中 $x\in A$，故
\[
  y=x+c\leqslant\alpha+c.
\]
所以 $\alpha+c$ 是 $B$ 的上界。另一方面，任取 $\varepsilon>0$，由上确界定义可取 $x\in A$ 使
\[
  x>\alpha-\varepsilon.
\]
于是 $x+c\in B$ 且
\[
  x+c>\alpha+c-\varepsilon.
\]
故任何小于 $\alpha+c$ 的数都不能作为 $B$ 的上界，所以
\[
  \sup B=\alpha+c=\sup A+c.
\]
对下确界，若 $\beta=\inf A$，则 $\beta+c$ 是 $B$ 的下界，并且对每个 $\varepsilon>0$ 可取 $x\in A$ 使 $x<\beta+\varepsilon$，从而 $x+c<\beta+c+\varepsilon$。故
\[
  \inf B=\inf A+c.
\]
上述结论在采用 $\pm\infty$ 约定时也按同样方式成立。
\end{xsolution}

\begin{xexercise}[3.1.3 练习题 8]
设 $C\subset\{x+y\mid x\in A,y\in B\}$，证明 $\sup C\leqslant\sup A+\sup B$，并举严格不等号的例子。
\end{xexercise}
\begin{xsolution}
记 $\alpha=\sup A$，$\beta=\sup B$。若 $c\in C$，则存在 $x\in A,y\in B$ 使 $c=x+y$，于是
\[
  c=x+y\leqslant\alpha+\beta.
\]
故 $\alpha+\beta$ 是 $C$ 的上界，因此
\[
  \sup C\leqslant\sup A+\sup B.
\]
严格不等号例如取 $A=B=(0,1)$，$C=\{1\}$。此时
\[
  \sup C=1<2=\sup A+\sup B.
\]
\end{xsolution}

\begin{xexercise}[3.1.3 练习题 9]
设 $C\supset\{x+y\mid x\in A,y\in B\}$，证明 $\sup C\geqslant\sup A+\sup B$，并举严格不等号的例子。
\end{xexercise}
\begin{xsolution}
仍记 $\alpha=\sup A$，$\beta=\sup B$。先证明
\[
  \sup\{x+y\mid x\in A,y\in B\}=\alpha+\beta.
\]
上界方向由上一题的论证立即得到。反之，任取 $\varepsilon>0$，可取 $x\in A,y\in B$ 使
\[
  x>\alpha-\frac\varepsilon2,
  \qquad
  y>\beta-\frac\varepsilon2.
\]
于是
\[
  x+y>\alpha+\beta-\varepsilon,
\]
故该和集的上确界不小于 $\alpha+\beta$。因此二者相等。

由于和集包含于 $C$，有
\[
  \sup C\geqslant\alpha+\beta.
\]
严格不等号例如取 $A=B=\{0\}$，$C=\{0,1\}$，则
\[
  \sup C=1>0=\sup A+\sup B.
\]
\end{xsolution}

\section{闭区间套定理}
\subsection*{3.2.3 练习题}

\begin{xexercise}[3.2.3 练习题 1]
对数列 $\{(-1)^n\}$ 从区间 $[-1,1]$ 出发，按例题 3.2.2 的方法构造闭区间套和收敛子列，并回答能否找出三个收敛子列。
\end{xexercise}
\begin{xsolution}
从 $I_1=[-1,1]$ 出发，二分后选择含有无穷多个数列项 $1$ 的右半区间，得到
\[
  I_2=[0,1],\quad I_3=[1/2,1],\quad I_4=[3/4,1],\ \ldots
\]
即对 $k\geqslant2$ 可取
\[
  I_k=[1-2^{2-k},1].
\]
每个 $I_k$ 都含有无穷多个偶数项 $x_{2m}=1$，而 $\abs{I_k}\to0$，交集为 $\{1\}$。在 $I_k$ 中依次选取下标递增的偶数项，例如取 $x_{2k}=1$，便得到收敛子列 $x_{2k}\to1$。

可以用同样方法得到三个不同的收敛子列。例如
\[
  \{x_{2k}\},\qquad \{x_{4k}\},\qquad \{x_{4k+2}\}
\]
都恒等于 $1$，因而都收敛于 $1$。也可在第一步选择左半区间构造收敛于 $-1$ 的奇数项子列。
\end{xsolution}

\begin{xexercise}[3.2.3 练习题 2]
把闭区间套改为开区间套而其余条件不变，举例说明结论不成立。
\end{xexercise}
\begin{xsolution}
取
\[
  I_n=(0,1/n),\qquad n\in\NN_+.
\]
则 $I_{n+1}\subset I_n$，并且 $\abs{I_n}=1/n\to0$，但
\[
  \bigcap_{n=1}^{\infty}(0,1/n)=\varnothing.
\]
因为若 $x$ 属于所有这些区间，则 $x>0$，但取 $n>1/x$ 就有 $1/n<x$，矛盾。
\end{xsolution}

\begin{xexercise}[3.2.3 练习题 3]
设 $\{(a_n,b_n)\}$ 为开区间套，$\{a_n\}$ 严格增加，$\{b_n\}$ 严格减少，且 $a_n<b_n$。证明其交非空。
\end{xexercise}
\begin{xsolution}
由包含关系，$\{a_n\}$ 单调增加并以任一 $b_m$ 为上界；$\{b_n\}$ 单调减少并以任一 $a_m$ 为下界。因此存在有限极限
\[
  \alpha=\lim_{n\to\infty}a_n,
  \qquad
  \beta=\lim_{n\to\infty}b_n,
\]
且由 $a_n<b_n$ 得 $\alpha\leqslant\beta$。

严格单调性给出对每个 $n$ 都有
\[
  a_n<\alpha,\qquad \beta<b_n.
\]
取
\[
  \xi=\frac{\alpha+\beta}{2}.
\]
则对每个 $n$，
\[
  a_n<\alpha\leqslant\xi\leqslant\beta<b_n,
\]
故 $\xi\in(a_n,b_n)$ 对所有 $n$ 成立。因此
\[
  \bigcap_{n=1}^{\infty}(a_n,b_n)\neq\varnothing.
\]
\end{xsolution}

\begin{xexercise}[3.2.3 练习题 4]
用闭区间套定理证明确界存在定理。
\end{xexercise}
\begin{xsolution}
只证有上界的非空数集 $A$ 有上确界。任取 $u_1\in A$。若 $u_1$ 已是 $A$ 的上界，则 $u_1=\max A=\sup A$。否则取 $v_1$ 为 $A$ 的一个上界，于是 $u_1<v_1$，并且 $u_1$ 不是上界、$v_1$ 是上界。

对区间 $[u_n,v_n]$ 二分，令 $c_n=(u_n+v_n)/2$。若 $c_n$ 是 $A$ 的上界，置
\[
  u_{n+1}=u_n,\qquad v_{n+1}=c_n;
\]
若 $c_n$ 不是 $A$ 的上界，置
\[
  u_{n+1}=c_n,\qquad v_{n+1}=v_n.
\]
于是 $\{[u_n,v_n]\}$ 是闭区间套，且始终保持
\[
  u_n\text{ 不是 }A\text{ 的上界},\qquad v_n\text{ 是 }A\text{ 的上界},
\]
并有
\[
  v_n-u_n=2^{1-n}(v_1-u_1)\to0.
\]
由闭区间套定理，存在唯一 $\xi$ 使 $u_n\to\xi$、$v_n\to\xi$。

任取 $x\in A$。因每个 $v_n$ 都是上界，$x\leqslant v_n$，令 $n\to\infty$ 得 $x\leqslant\xi$，故 $\xi$ 是 $A$ 的上界。另一方面，若 $y<\xi$，则因 $u_n\to\xi$，可取 $n$ 使 $u_n>y$。由于 $u_n$ 不是 $A$ 的上界，存在 $x\in A$ 使 $x>u_n>y$，故 $y$ 不是上界。因此 $\xi$ 是最小上界，即 $\xi=\sup A$。

对下确界应用同样构造，或对集合 $-A$ 使用上确界结论，即得下确界存在定理。
\end{xsolution}

\begin{xexercise}[3.2.3 练习题 5]
用闭区间套定理证明单调有界数列的收敛定理。
\end{xexercise}
\begin{xsolution}
先设 $\{a_n\}$ 单调增加且有上界。由上一题，闭区间套定理推出确界存在定理，因此数集
\[
  A=\{a_n\mid n\in\NN_+\}
\]
有上确界 $L=\sup A$。任取 $\varepsilon>0$，$L-\varepsilon$ 不是 $A$ 的上界，故存在 $N$ 使
\[
  a_N>L-\varepsilon.
\]
由单调性，当 $n\geqslant N$ 时
\[
  L-\varepsilon<a_N\leqslant a_n\leqslant L,
\]
于是 $\abs{a_n-L}<\varepsilon$，即 $a_n\to L$。

若 $\{a_n\}$ 单调减少且有下界，则对 $-a_n$ 应用上面的结论，或用下确界作同样证明，即知其收敛。
\end{xsolution}

\section{凝聚定理}
\subsection*{3.3.1 思考题}

\begin{xthought}[3.3.1 思考题]
说明凝聚定理在有理数集 $\QQ$ 中不成立。
\end{xthought}
\begin{xsolution}
取一列有理数 $r_n$，使 $r_n\to\sqrt2$，例如取 $\sqrt2$ 的有限十进不足近似
\[
  1,\ 1.4,\ 1.41,\ 1.414,\ \ldots.
\]
这是 $\QQ$ 中的有界数列。若它在 $\QQ$ 中有收敛子列 $r_{n_k}\to q\in\QQ$，由于原数列在 $\RR$ 中收敛于 $\sqrt2$，它的每个子列在 $\RR$ 中也收敛于 $\sqrt2$。极限唯一性将给出 $q=\sqrt2$，与 $q\in\QQ$ 矛盾。因此这个有界有理数列在 $\QQ$ 中没有收敛子列，凝聚定理在 $\QQ$ 中不成立。
\end{xsolution}

\subsection*{3.3.3 练习题}

\begin{xexercise}[3.3.3 练习题 1]
证明：在 $a$ 的每个邻域中有数列 $\{x_n\}$ 的无穷多项，当且仅当 $a$ 是某个子列的极限。
\end{xexercise}
\begin{xsolution}
先证必要性。设 $a$ 的每个邻域都含有 $\{x_n\}$ 的无穷多项。对 $k=1,2,\ldots$，在邻域
\[
  \left(a-\frac1k,a+\frac1k\right)
\]
中递归选取一项 $x_{n_k}$，并要求 $n_k>n_{k-1}$。这是可能的，因为该邻域含有无穷多项。于是
\[
  \abs{x_{n_k}-a}<\frac1k,
\]
故 $x_{n_k}\to a$。

反之，若存在子列 $x_{n_k}\to a$，则对 $a$ 的任一邻域 $U$，从某个 $K$ 起所有 $x_{n_k}$ 都属于 $U$。因而 $U$ 中含有原数列的无穷多项。
\end{xsolution}

\begin{xexercise}[3.3.3 练习题 2]
证明：有界数列发散的充分必要条件是存在两个收敛于不同极限值的子列。
\end{xexercise}
\begin{xsolution}
若存在两个子列分别收敛于 $a\neq b$，则原数列不可能收敛，因为收敛数列的每个子列都收敛于原极限。

反之，设 $\{x_n\}$ 有界但发散。由凝聚定理，它有一个收敛子列 $x_{n_k}\to a$。若不存在收敛于不同极限的子列，则所有收敛子列的极限都只能是 $a$。我们证明这将迫使 $x_n\to a$。

若 $x_n$ 不收敛于 $a$，则存在 $\varepsilon_0>0$，使得对任意 $N$ 都能找到 $n>N$ 满足
\[
  \abs{x_n-a}\geqslant\varepsilon_0.
\]
从这些项中选出子列 $\{x_{m_j}\}$。它仍然有界，故由凝聚定理又有收敛子列 $x_{m_{j_l}}\to b$。对不等式取极限得
\[
  \abs{b-a}\geqslant\varepsilon_0,
\]
所以 $b\neq a$，矛盾。因此必存在两个收敛于不同极限值的子列。
\end{xsolution}

\begin{xexercise}[3.3.3 练习题 3]
证明：若 $\{x_n\}$ 无界，但不是无穷大量，则存在一个收敛子列和一个无穷大量子列。
\end{xexercise}
\begin{xsolution}
因为 $\{x_n\}$ 不是无穷大量，所以“对每个 $G>0$ 最终都有 $\abs{x_n}>G$”不成立。故存在 $G_0>0$，使得对任意 $N$ 都有某个 $n>N$ 满足 $\abs{x_n}\leqslant G_0$。递归选取这些项，得到一个有界子列。由凝聚定理，该有界子列有收敛子列。

另一方面，原数列无界。递归选取 $n_k>n_{k-1}$ 使
\[
  \abs{x_{n_k}}>k.
\]
于是 $\abs{x_{n_k}}\to+\infty$，故 $\{x_{n_k}\}$ 是无穷大量子列。
\end{xsolution}

\begin{xexercise}[3.3.3 练习题 4]
用凝聚定理证明单调有界数列的收敛定理。
\end{xexercise}
\begin{xsolution}
只讨论单调增加且有上界的数列 $\{a_n\}$。由凝聚定理，存在子列 $a_{n_k}\to L$。先说明对每个 $n$ 都有 $a_n\leqslant L$：固定 $n$，取充分大的 $k$ 使 $n_k\geqslant n$，则
\[
  a_n\leqslant a_{n_k}.
\]
令 $k\to\infty$ 得 $a_n\leqslant L$。

任取 $\varepsilon>0$。由 $a_{n_k}\to L$，可取 $K$ 使
\[
  L-\varepsilon<a_{n_K}\leqslant L.
\]
当 $n\geqslant n_K$ 时，由单调性
\[
  L-\varepsilon<a_{n_K}\leqslant a_n\leqslant L,
\]
故 $\abs{a_n-L}<\varepsilon$。于是 $a_n\to L$。若 $\{a_n\}$ 单调减少且有下界，则 $\{-a_n\}$ 单调增加且有上界，由刚证结论存在 $M$ 使 $-a_n\to M$，于是 $a_n\to-M$。
\end{xsolution}

\section{Cauchy 收敛准则}
\subsection*{3.4.1 思考题}

\begin{xthought}[3.4.1 思考题]
说明 Cauchy 收敛准则在有理数集 $\QQ$ 中不成立。
\end{xthought}
\begin{xsolution}
仍取有理数列 $r_n\to\sqrt2$。在 $\RR$ 中它是收敛数列，因此是 Cauchy 数列；Cauchy 性只涉及各项之间的距离，所以把它看作 $\QQ$ 中的数列时仍是 Cauchy 数列。但它在 $\QQ$ 中没有极限，因为若 $r_n\to q\in\QQ$，由实数极限唯一性必有 $q=\sqrt2\notin\QQ$。因此在 $\QQ$ 中“基本数列必收敛”失败。
\end{xsolution}

\subsection*{3.4.5 练习题}

\begin{xexercise}[3.4.5 练习题 1]
判断原题四个条件能否保证 $\{x_n\}$ 是基本数列。
\end{xexercise}
\begin{xsolution}
\begin{enumerate}[label=(\arabic*)]
  \item 能。给定 $\varepsilon>0$，把题设中的误差取为 $\varepsilon/2$，则存在 $N$，当 $n>N$ 时
  \[
    \abs{x_n-x_N}<\frac\varepsilon2.
  \]
  因而对任意 $m,n>N$，
  \[
    \abs{x_n-x_m}
    \leqslant\abs{x_n-x_N}+\abs{x_m-x_N}<\varepsilon.
  \]
  所以 $\{x_n\}$ 是基本数列。

  \item 不能。取调和部分和
  \[
    x_n=1+\frac12+\cdots+\frac1n.
  \]
  对任意 $p\geqslant1$，
  \[
    \abs{x_{n+p}-x_n}
    =\sum_{k=n+1}^{n+p}\frac1k
    \leqslant\frac{p}{n}<\frac{p}{n}+0,
  \]
  因而满足题设估计，但调和部分和发散，所以不是基本数列。

  \item 能。取 $p=1$ 得
  \[
    \abs{x_{n+1}-x_n}\leqslant\frac1{n^2}.
  \]
  若 $m>n\geqslant2$，则
  \[
    \abs{x_m-x_n}
    \leqslant\sum_{k=n}^{m-1}\frac1{k^2}
    \leqslant\sum_{k=n}^{\infty}\frac1{k(k-1)}
    =\frac1{n-1}.
  \]
  右端趋于 $0$，故 $\{x_n\}$ 是基本数列。

  \item 不能。仍取调和部分和 $x_n=1+1/2+\cdots+1/n$。对每个固定的正整数 $p$，
  \[
    \abs{x_{n+p}-x_n}
    =\sum_{k=n+1}^{n+p}\frac1k
    \leqslant\frac pn\to0,
  \]
  但 $\{x_n\}$ 发散，因此不是基本数列。
\end{enumerate}
\end{xsolution}

\begin{xexercise}[3.4.5 练习题 2]
用对偶法则写出数列发散的 Cauchy 型充分必要条件。
\end{xexercise}
\begin{xsolution}
Cauchy 收敛准则为
\[
  \forall\varepsilon>0,\ \exists N,\ \forall m,n>N,
  \quad \abs{x_n-x_m}<\varepsilon.
\]
逐层否定量词，得到数列发散的充分必要条件：
\[
  \boxed{
  \exists\varepsilon_0>0,\ \forall N,\ \exists m,n>N,
  \quad \abs{x_n-x_m}\geqslant\varepsilon_0.}
\]
这里 $m,n$ 可随 $N$ 改变。
\end{xsolution}

\begin{xexercise}[3.4.5 练习题 3]
证明原题所列三个数列均为基本数列。
\end{xexercise}
\begin{xsolution}
\begin{enumerate}[label=(\arabic*)]
  \item 对 $m>n$，
  \[
    \abs{a_m-a_n}
    =\sum_{k=n+1}^{m}\frac1{k!}.
  \]
  当 $k\geqslant n+1$ 时，后继项与前项之比至多为 $1/(n+2)\leqslant1/2$，故
  \[
    \abs{a_m-a_n}
    \leqslant\frac1{(n+1)!}\sum_{j=0}^{\infty}\frac1{2^j}
    =\frac2{(n+1)!}\to0.
  \]
  因而 $\{a_n\}$ 是基本数列。

  \item 对交错调和部分和 $b_n$，若 $m>n$，交错和余段的绝对值不超过首项，故
  \[
    \abs{b_m-b_n}\leqslant\frac1{n+1}\to0.
  \]
  因而 $\{b_n\}$ 是基本数列。这里的估计可直接由把相邻两项配对得到。

  \item 设
  \[
    u_k=\frac{\sin kx}{k(k+\sin kx)},\qquad k\geqslant2.
  \]
  因 $k+\sin kx\geqslant k-1>0$，有
  \[
    \abs{u_k}\leqslant\frac1{k(k-1)}
    =\frac1{k-1}-\frac1k.
  \]
  故 $m>n\geqslant2$ 时
  \[
    \abs{c_m-c_n}
    \leqslant\sum_{k=n+1}^{m}\frac1{k(k-1)}
    =\frac1n-\frac1m<\frac1n\to0.
  \]
  因此 $\{c_n\}$ 是基本数列。
\end{enumerate}
\end{xsolution}

\begin{xexercise}[3.4.5 练习题 4]
设 $a_n=\sin1+\frac{\sin2}{2!}+\cdots+\frac{\sin n}{n!}$。证明：(1) 有界但不单调；(2) 收敛。
\end{xexercise}
\begin{xsolution}
\begin{enumerate}[label=(\arabic*)]
  \item 由 $\abs{\sin k}\leqslant1$，
  \[
    \abs{a_n}\leqslant\sum_{k=1}^{n}\frac1{k!}
    <\sum_{k=1}^{\infty}\frac1{k!}<\infty,
  \]
  故数列有界。又
  \[
    a_2-a_1=\frac{\sin2}{2!}>0,
    \qquad
    a_4-a_3=\frac{\sin4}{4!}<0,
  \]
  因为 $0<2<\pi$ 而 $\pi<4<2\pi$。所以它既非单调增加也非单调减少。

  \item 若 $m>n$，则
  \[
    \abs{a_m-a_n}
    \leqslant\sum_{k=n+1}^{m}\frac1{k!}
    \leqslant\frac2{(n+1)!}\to0.
  \]
  因此 $\{a_n\}$ 是基本数列，由 Cauchy 收敛准则知它收敛。
\end{enumerate}
\end{xsolution}

\begin{xexercise}[3.4.5 练习题 5]
若 $x_n=\sum_{k=1}^{n}a_k$，$y_n=\sum_{k=1}^{n}\abs{a_k}$，且 $\{y_n\}$ 收敛，证明 $\{x_n\}$ 收敛。
\end{xexercise}
\begin{xsolution}
由 $y_n$ 收敛，$\{y_n\}$ 是基本数列。给定 $\varepsilon>0$，存在 $N$，当 $m>n>N$ 时
\[
  y_m-y_n=\sum_{k=n+1}^{m}\abs{a_k}<\varepsilon.
\]
于是
\[
  \abs{x_m-x_n}
  =\abs{\sum_{k=n+1}^{m}a_k}
  \leqslant\sum_{k=n+1}^{m}\abs{a_k}
  =y_m-y_n<\varepsilon.
\]
故 $\{x_n\}$ 是基本数列，因此收敛。
\end{xsolution}

\begin{xexercise}[3.4.5 练习题 6]
设 $S_n=1+2^{-p}+3^{-p}+\cdots+n^{-p}$，其中 $p\leqslant1$。证明 $\{S_n\}$ 发散。
\end{xexercise}
\begin{xsolution}
对任意 $n$，有
\[
  S_{2n}-S_n=\sum_{k=n+1}^{2n}\frac1{k^p}.
\]
若 $p\leqslant0$，则 $k^{-p}\geqslant1$，所以
\[
  S_{2n}-S_n=\sum_{k=n+1}^{2n}k^{-p}\geqslant n,
\]
于是任取 $N$，选 $n>N$ 并令 $m=2n$，都有 $\abs{S_m-S_n}\geqslant n>1$，故取 $\varepsilon_0=1$ 即知 Cauchy 条件不成立。

若 $0<p\leqslant1$，则对 $n<k\leqslant2n$ 有
\[
  \frac1{k^p}\geqslant\frac1{(2n)^p},
\]
因此
\[
  S_{2n}-S_n
  \geqslant\frac{n}{(2n)^p}
  =2^{-p}n^{1-p}\geqslant2^{-p}>0.
\]
取固定的 $\varepsilon_0=2^{-p}/2$，无论 $N$ 多大，都可选 $n>N$，令 $m=2n$，使
\[
  \abs{S_m-S_n}>\varepsilon_0.
\]
故 $\{S_n\}$ 不是基本数列，从而发散。
\end{xsolution}

\begin{xexercise}[3.4.5 练习题 7]
对 $0<q<1$，证明迭代 $x_{n+1}=q\sin x_n+a$ 求解 Kepler 方程 $x-q\sin x=a$ 的方法正确。
\end{xexercise}
\begin{xsolution}
令
\[
  f(x)=q\sin x+a.
\]
利用 $\abs{\sin x-\sin y}\leqslant\abs{x-y}$，有
\[
  \abs{f(x)-f(y)}\leqslant q\abs{x-y},
\]
其中 $0<q<1$。于是
\[
  \abs{x_{n+1}-x_n}
  \leqslant q\abs{x_n-x_{n-1}}
  \leqslant q^{n-1}\abs{x_2-x_1}.
\]
若 $p\geqslant1$，则
\[
  \abs{x_{n+p}-x_n}
  \leqslant\sum_{j=n}^{n+p-1}\abs{x_{j+1}-x_j}
  \leqslant\frac{q^{n-1}}{1-q}\abs{x_2-x_1}\to0.
\]
所以 $\{x_n\}$ 是基本数列，存在极限 $\xi$。

又由压缩估计
\[
  \abs{f(x_n)-f(\xi)}\leqslant q\abs{x_n-\xi}\to0,
\]
故令 $n\to\infty$ 于 $x_{n+1}=f(x_n)$ 得
\[
  \xi=q\sin\xi+a,
\]
即 $\xi$ 是 Kepler 方程的解。

若 $\eta$ 也是解，则
\[
  \abs{\xi-\eta}
  =q\abs{\sin\xi-\sin\eta}
  \leqslant q\abs{\xi-\eta}.
\]
因 $q<1$，必有 $\xi=\eta$。所以解唯一，且从任意初值开始的上述迭代都收敛到这个唯一解，方法正确。
\end{xsolution}

\section{覆盖定理}
\subsection*{3.5.3 练习题}

\begin{xexercise}[3.5.3 练习题 1]
对开区间 $(0,1)$ 构造一个开覆盖，使它的每一个有限子集都不能覆盖 $(0,1)$。
\end{xexercise}
\begin{xsolution}
取
\[
  O_n=(1/n,1),\qquad n=2,3,\ldots.
\]
任意 $x\in(0,1)$，取 $n>1/x$，则 $1/n<x<1$，故
\[
  (0,1)=\bigcup_{n=2}^{\infty}O_n.
\]
但任取有限多个 $O_{n_1},\ldots,O_{n_k}$，令 $N=\max\{n_1,\ldots,n_k\}$，则
\[
  \bigcup_{j=1}^{k}O_{n_j}=(1/N,1),
\]
它漏掉了 $(0,1/N]$ 中的点。因此没有有限子覆盖。
\end{xsolution}

\begin{xexercise}[3.5.3 练习题 2]
用闭区间套定理证明覆盖定理。
\end{xexercise}
\begin{xsolution}
设闭区间 $[a,b]$ 有开覆盖 $\{O_\alpha\}$。反设不存在有限子覆盖。二分 $[a,b]$，若两个半区间都有有限子覆盖，把两组有限子覆盖合并便得到 $[a,b]$ 的有限子覆盖，矛盾。因此至少有一个半区间没有有限子覆盖，把它记为 $I_1$。

继续二分 $I_1$。若它的两个半区间都有有限子覆盖，则合并这两个有限子覆盖就会覆盖 $I_1$，与 $I_1$ 的选取矛盾。因此仍至少有一个半区间没有有限子覆盖，记为 $I_2$。归纳得到闭区间套
\[
  I_1\supset I_2\supset\cdots,
\]
其中每个 $I_n$ 都没有有限子覆盖，而且
\[
  \abs{I_n}=2^{-n}(b-a)\to0.
\]
由闭区间套定理，存在唯一 $\xi\in\bigcap_{n=1}^{\infty}I_n$。

因为 $\{O_\alpha\}$ 覆盖 $[a,b]$，存在某个 $O_{\alpha_0}$ 使 $\xi\in O_{\alpha_0}$。因 $O_{\alpha_0}$ 是开区间，可取 $\delta>0$ 使
\[
  (\xi-\delta,\xi+\delta)\subset O_{\alpha_0}.
\]
取 $n$ 足够大使 $\abs{I_n}<\delta$。由于 $\xi\in I_n$，整个 $I_n$ 都包含在 $(\xi-\delta,\xi+\delta)$ 中，因而单个 $O_{\alpha_0}$ 就覆盖 $I_n$。这与 $I_n$ 没有有限子覆盖矛盾。因此 $[a,b]$ 必有有限子覆盖。
\end{xsolution}

\begin{xexercise}[3.5.3 练习题 3]
用覆盖定理证明闭区间套定理。
\end{xexercise}
\begin{xsolution}
设 $I_n=[a_n,b_n]$ 为闭区间套。先证交集非空。反设
\[
  \bigcap_{n=1}^{\infty}I_n=\varnothing.
\]
于是对每个 $x\in I_1$，存在 $n$ 使 $x\notin I_n$，即
\[
  x<a_n\quad\text{或}\quad x>b_n.
\]
因此开区间族
\[
  \{(-\infty,a_n),(b_n,+\infty)\mid n\in\NN_+\}
\]
覆盖 $I_1$。由覆盖定理可取有限子覆盖。设其中出现的最大下标为 $N$。由于 $a_n$ 单调不减、$b_n$ 单调不增，这个有限子覆盖的并包含于
\[
  (-\infty,a_N)\cup(b_N,+\infty),
\]
从而不可能覆盖非空闭区间 $I_N=[a_N,b_N]\subset I_1$，矛盾。因此交集非空。

若再有 $\abs{I_n}=b_n-a_n\to0$，假设交集中有两个不同点 $x<y$，则对每个 $n$ 都有 $x,y\in I_n$，所以
\[
  b_n-a_n\geqslant y-x>0,
\]
与长度趋于 $0$ 矛盾。故交集只能是单点集。
\end{xsolution}

\begin{xexercise}[3.5.3 练习题 4]
用覆盖定理证明凝聚定理。
\end{xexercise}
\begin{xsolution}
设有界数列 $\{x_n\}$ 的所有项都在某个闭区间 $[a,b]$ 中。反设它没有收敛子列。于是对每个 $\xi\in[a,b]$，必存在 $\xi$ 的一个开邻域 $O(\xi)$，其中只含数列的有限多项；否则若 $\xi$ 的每个邻域都含无穷多项，就可依次在半径 $1/k$ 的邻域中选取下标递增的 $x_{n_k}$，得到 $x_{n_k}\to\xi$。

于是 $\{O(\xi)\mid\xi\in[a,b]\}$ 是 $[a,b]$ 的开覆盖。由覆盖定理，有有限子覆盖
\[
  O(\xi_1),\ldots,O(\xi_m).
\]
每个邻域中只含原数列的有限多项，因此这有限个邻域的并也只能含有限多项；但它们覆盖 $[a,b]$，而原数列的所有项都在 $[a,b]$ 中，矛盾。故有界数列必有收敛子列。
\end{xsolution}

\begin{xexercise}[3.5.3 练习题 5]
给例题 3.5.2 中的数集 $A$ 举两个具体例子：(1) $A$ 无上界；(2) $A$ 有上界，且 $b<\xi=\sup A$、$\xi\notin A$。
\end{xexercise}
\begin{xsolution}
在例题 3.5.2 中，固定左端点 $a$，定义
\[
  A=\{x\geqslant a\mid [a,x]\text{ 有有限子覆盖}\}.
\]
可取以下具体例子。
\begin{enumerate}[label=(\arabic*)]
  \item 令 $a=0,b=1$，开覆盖族取
  \[
    O_m=(-1,m),\qquad m\in\NN_+.
  \]
  它当然覆盖 $[0,1]$。对任意 $x\geqslant0$，选整数 $m>x$，则单个 $O_m$ 已覆盖 $[0,x]$，所以
  \[
    A=[0,+\infty),
  \]
  无上界。

  \item 仍令 $a=0,b=1$，只取一个开区间
  \[
    O=(-1,2).
  \]
  它覆盖 $[0,1]$。这时 $[0,x]$ 能被该开覆盖有限覆盖，当且仅当 $0\leqslant x<2$。因此
  \[
    A=[0,2),\qquad \xi=\sup A=2>b=1,
  \]
  且 $2\notin A$。
\end{enumerate}
\end{xsolution}

\section{数列的上极限和下极限}
\subsection*{3.6.4 练习题}

\begin{xexercise}[3.6.4 练习题 1]
求原题所列四个数列的上极限和下极限。
\end{xexercise}
\begin{xsolution}
\begin{enumerate}[label=(\arabic*)]
  \item $x_n=(1+(-1)^n)/2$ 的奇数项恒为 $0$、偶数项恒为 $1$，故
  \[
    \underline{\lim}_{n\to\infty}x_n=0,
    \qquad
    \overline{\lim}_{n\to\infty}x_n=1.
  \]

  \item $x_n=\sin(n\pi/4)$ 以周期 $8$ 重复取值
  \[
    \frac{\sqrt2}{2},\ 1,\ \frac{\sqrt2}{2},\ 0,
    -\frac{\sqrt2}{2},\ -1,\ -\frac{\sqrt2}{2},\ 0.
  \]
  因此
  \[
    \underline{\lim}_{n\to\infty}x_n=-1,
    \qquad
    \overline{\lim}_{n\to\infty}x_n=1.
  \]

  \item 对 $x_n=n^{(-1)^n}$，偶数项 $x_{2k}=2k\to+\infty$，奇数项 $x_{2k-1}=1/(2k-1)\to0$。故
  \[
    \underline{\lim}_{n\to\infty}x_n=0,
    \qquad
    \overline{\lim}_{n\to\infty}x_n=+\infty.
  \]

  \item 对 $x_n=\ee^{n(-1)^n}$，偶数项趋于 $+\infty$，奇数项 $\ee^{-n}\to0$。故
  \[
    \underline{\lim}_{n\to\infty}x_n=0,
    \qquad
    \overline{\lim}_{n\to\infty}x_n=+\infty.
  \]
\end{enumerate}
\end{xsolution}

\begin{xexercise}[3.6.4 练习题 2]
若 $x_n\geqslant y_n$，证明上下极限都保持这一不等式方向。
\end{xexercise}
\begin{xsolution}
对每个 $n$，由 $x_k\geqslant y_k$（$k\geqslant n$）有
\[
  \sup_{k\geqslant n}x_k\geqslant\sup_{k\geqslant n}y_k,
  \qquad
  \inf_{k\geqslant n}x_k\geqslant\inf_{k\geqslant n}y_k.
\]
令 $n\to\infty$，利用上下极限的尾集表示，得到
\[
  \overline{\lim}_{n\to\infty}x_n
  \geqslant\overline{\lim}_{n\to\infty}y_n,
  \qquad
  \underline{\lim}_{n\to\infty}x_n
  \geqslant\underline{\lim}_{n\to\infty}y_n.
\]
\end{xsolution}

\begin{xexercise}[3.6.4 练习题 3]
令 $S_n=\frac1n\sum_{k=1}^{n}x_k$。证明
\[
\underline{\lim}x_n\leqslant\underline{\lim}S_n
\leqslant\overline{\lim}S_n\leqslant\overline{\lim}x_n.
\]
\end{xexercise}
\begin{xsolution}
记
\[
  l=\underline{\lim}_{n\to\infty}x_n,
  \qquad
  L=\overline{\lim}_{n\to\infty}x_n.
\]
先设 $l,L$ 有限。任取 $\varepsilon>0$。由上下极限的定义，存在 $N$，当 $k\geqslant N$ 时
\[
  l-\varepsilon<x_k<L+\varepsilon.
\]
于是对 $n>N$，
\[
  \frac1n\sum_{k=1}^{N-1}x_k
  +\frac{n-N+1}{n}(l-\varepsilon)
  <S_n
  <\frac1n\sum_{k=1}^{N-1}x_k
  +\frac{n-N+1}{n}(L+\varepsilon).
\]
令 $n\to\infty$，有限个初始项的平均趋于 $0$，得到
\[
  l-\varepsilon\leqslant\underline{\lim}S_n
  \leqslant\overline{\lim}S_n\leqslant L+\varepsilon.
\]
再令 $\varepsilon\downarrow0$，即得所需不等式。

若 $l=-\infty$ 或 $L=+\infty$，相应外侧不等式自动成立；若只有一侧为有限数，则上述相应的单边估计仍然有效。因此扩展实数情形也成立。
\end{xsolution}

\begin{xexercise}[3.6.4 练习题 4]
设 $\{a_n\}$ 为正数列。证明
\[
  \overline{\lim}_{n\to\infty}\sqrt[n]{a_n}\leqslant1
\]
当且仅当对每个 $l>1$ 都有 $a_n/l^n\to0$。
\end{xexercise}
\begin{xsolution}
先设
\[
  \overline{\lim}_{n\to\infty}\sqrt[n]{a_n}\leqslant1.
\]
任取 $l>1$，再取 $r$ 满足 $1<r<l$。由上极限定义，存在 $N$，当 $n>N$ 时
\[
  \sqrt[n]{a_n}<r.
\]
于是
\[
  0<\frac{a_n}{l^n}<\left(\frac rl\right)^n\to0,
\]
故 $a_n/l^n\to0$。

反之，假设对每个 $l>1$ 都有 $a_n/l^n\to0$。若
\[
  \overline{\lim}_{n\to\infty}\sqrt[n]{a_n}>1,
\]
则可取 $1<l<r<\overline{\lim}\sqrt[n]{a_n}$。由上极限的定义，有无穷多个 $n$ 满足 $\sqrt[n]{a_n}>r$，从而这些 $n$ 上
\[
  \frac{a_n}{l^n}>\left(\frac rl\right)^n.
\]
右端沿该子列趋于 $+\infty$，与 $a_n/l^n\to0$ 矛盾。因此上极限不超过 $1$。
\end{xsolution}

\begin{xexercise}[3.6.4 练习题 5]
设 $\{a_n\}$ 为正数列。证明比值的上下极限夹住 $\sqrt[n]{a_n}$ 的上下极限。
\end{xexercise}
\begin{xsolution}
记
\[
  r=\underline{\lim}_{n\to\infty}\frac{a_{n+1}}{a_n},
  \qquad
  R=\overline{\lim}_{n\to\infty}\frac{a_{n+1}}{a_n}.
\]
先证右端不等式。若 $R<+\infty$，任取 $q>R$。存在 $N$，当 $n\geqslant N$ 时
\[
  \frac{a_{n+1}}{a_n}<q.
\]
故对 $n>N$，
\[
  a_n\leqslant a_N q^{n-N}.
\]
开 $n$ 次方得
\[
  \sqrt[n]{a_n}\leqslant a_N^{1/n}q^{1-N/n}\to q,
\]
所以
\[
  \overline{\lim}_{n\to\infty}\sqrt[n]{a_n}\leqslant q.
\]
令 $q\downarrow R$，得到
\[
  \overline{\lim}\sqrt[n]{a_n}\leqslant R.
\]
若 $R=+\infty$，此式自动成立。

再证左端不等式。若 $r>0$，任取 $0<q<r$。充分大的 $n$ 有
\[
  \frac{a_{n+1}}{a_n}>q,
\]
从而 $a_n\geqslant a_Nq^{n-N}$，于是
\[
  \underline{\lim}_{n\to\infty}\sqrt[n]{a_n}\geqslant q.
\]
令 $q\uparrow r$ 得
\[
  r\leqslant\underline{\lim}\sqrt[n]{a_n}.
\]
若 $r=0$，因根式非负，此式也自动成立。中间
\[
  \underline{\lim}\sqrt[n]{a_n}\leqslant\overline{\lim}\sqrt[n]{a_n}
\]
由定义成立。合并即得原题全部不等式。
\end{xsolution}

\begin{xexercise}[3.6.4 练习题 6]
根据极限点定义直接证明：每个数列的极限点集中必有最大数和最小数。
\end{xexercise}
\begin{xsolution}
先证最大极限点存在，不使用命题 3.6.2、3.6.3。

若 $\{x_n\}$ 无上界，则可递归选取 $n_k$ 使 $x_{n_k}>k$，所以 $x_{n_k}\to+\infty$；按本章约定 $+\infty$ 是极限点，并且在扩展实数的次序中没有比它更大的点，因此它是最大极限点。以下设 $\{x_n\}$ 有上界。

设有限极限点集合为 $E$。若 $E\neq\varnothing$，则 $E$ 有上界。令 $L=\sup E$。对每个 $k$，可取 $c_k\in E$ 使
\[
  L-\frac1k<c_k\leqslant L.
\]
由于 $c_k$ 是极限点，在 $(c_k-1/k,c_k+1/k)$ 中有原数列的无穷多项。于是可递归选取 $n_k>n_{k-1}$ 使
\[
  \abs{x_{n_k}-c_k}<\frac1k.
\]
因此
\[
  \abs{x_{n_k}-L}
  \leqslant\abs{x_{n_k}-c_k}+\abs{c_k-L}<\frac2k,
\]
故 $x_{n_k}\to L$。所以 $L\in E$，即 $L$ 是最大有限极限点。

若 $E=\varnothing$，则数列必须是负无穷大量。否则存在实数 $M$，使得有无穷多项满足 $x_n\geqslant M$；结合数列有上界，这些项组成有界子列，由凝聚定理应有有限极限点，与 $E=\varnothing$ 矛盾。因此此时唯一的极限点是 $-\infty$，它同时也是最大极限点。

最小极限点的存在性完全对称：若数列无下界，则 $-\infty$ 为最小极限点；若有下界且有限极限点非空，对其取下确界并作同样的对角选取；若没有有限极限点，则数列为正无穷大量，唯一极限点为 $+\infty$。故最大、最小极限点都存在。
\end{xsolution}

\begin{xexercise}[3.6.4 练习题 7]
证明例题 3.6.2 中关于上、下极限所用的公式。
\end{xexercise}
\begin{xsolution}
例题中 $b_n\in[1.5,2]$，故所有量均为正且有限。记
\[
  \alpha=\underline{\lim}_{n\to\infty}b_n,
  \qquad
  \beta=\overline{\lim}_{n\to\infty}b_n.
\]
删除有限个初始项不改变上下极限，因此
\[
  \underline{\lim}b_{n+1}=\alpha,
  \qquad
  \overline{\lim}b_{n+1}=\beta.
\]
对尾集利用倒数函数在正数上的反序性，有
\[
  \sup_{k\geqslant n}\frac1{b_k}
  =\frac1{\inf_{k\geqslant n}b_k},
  \qquad
  \inf_{k\geqslant n}\frac1{b_k}
  =\frac1{\sup_{k\geqslant n}b_k}.
\]
令 $n\to\infty$，得到
\[
  \overline{\lim}_{n\to\infty}\frac1{b_n}=\frac1\alpha,
  \qquad
  \underline{\lim}_{n\to\infty}\frac1{b_n}=\frac1\beta.
\]
又常数相加可直接移出上下极限。于是由
\[
  b_{n+1}=1+\frac1{b_n}
\]
两边取下极限和上极限，分别得到
\[
  \alpha=1+\frac1\beta,
  \qquad
  \beta=1+\frac1\alpha,
\]
这正是例题中使用的公式。
\end{xsolution}

\begin{xexercise}[3.6.4 练习题 8]
证明公式 (3.2)：
\[
  \overline{\lim}_{n\to\infty}x_n
  =\lim_{n\to\infty}\sup_{k\geqslant n}\{x_k\}.
\]
\end{xexercise}
\begin{xsolution}
对数列 $\{-x_n\}$ 应用已经证明的公式 (3.3)，有
\[
  \underline{\lim}_{n\to\infty}(-x_n)
  =\lim_{n\to\infty}\inf_{k\geqslant n}(-x_k).
\]
利用
\[
  \overline{\lim}x_n=-\underline{\lim}(-x_n),
  \qquad
  \inf_{k\geqslant n}(-x_k)
  =-\sup_{k\geqslant n}x_k,
\]
两边同乘 $-1$，即得
\[
  \overline{\lim}_{n\to\infty}x_n
  =\lim_{n\to\infty}\sup_{k\geqslant n}x_k.
\]
扩展实数 $\pm\infty$ 的情形也由同一恒等式成立。
\end{xsolution}

\begin{xexercise}[3.6.4 练习题 9]
写出并证明与命题 3.6.3 对应的上极限判别。
\end{xexercise}
\begin{xsolution}
对应结论如下。

(1) 有限数 $b$ 是 $\{x_n\}$ 的上极限，当且仅当对每个 $\varepsilon>0$，邻域 $(b-\varepsilon,b+\varepsilon)$ 中有数列的无穷多项，并且存在 $N$，当 $n>N$ 时
\[
  x_n<b+\varepsilon.
\]

(2) $+\infty$ 是上极限，当且仅当 $\{x_n\}$ 无上界。

(3) $-\infty$ 是上极限，当且仅当 $\{x_n\}$ 为负无穷大量。

证明有限情形。若 $b$ 是上极限，则它是极限点，所以每个邻域有无穷多项。若对某个 $\varepsilon>0$ 有无穷多项满足 $x_n\geqslant b+\varepsilon$，从这些项中取子列；若该子列有界，凝聚定理给出一个不小于 $b+\varepsilon$ 的有限极限点；若无界，则有 $+\infty$ 极限点。两者都与 $b$ 为最大极限点矛盾。因此最终必有 $x_n<b+\varepsilon$。

反之，假设上述两条件成立。由每个邻域中有无穷多项，可递归选取 $n_k$ 使
\[
  \abs{x_{n_k}-b}<\frac1k,
\]
故 $b$ 是极限点。任取 $b'>b$，令 $\varepsilon=(b'-b)/2$。最终有
\[
  x_n<b+\varepsilon=b'-\varepsilon<b',
\]
所以不可能有子列收敛到 $b'$，也不可能有趋于 $+\infty$ 的子列。因此不存在大于 $b$ 的极限点，$b$ 即最大极限点。

无穷情形 (2)、(3) 可由极限点定义直接得到，或对命题 3.6.3 应用于 $-x_n$ 得到。
\end{xsolution}

\begin{xexercise}[3.6.4 练习题 10]
证明公式 (3.5)：
\[
  \overline{\lim}(x_n+y_n)
  \leqslant\overline{\lim}x_n+\overline{\lim}y_n.
\]
\end{xexercise}
\begin{xsolution}
设右端有意义。对每个 $n$，记
\[
  A_n=\sup_{k\geqslant n}x_k,
  \qquad
  B_n=\sup_{k\geqslant n}y_k.
\]
对所有 $k\geqslant n$ 有 $x_k\leqslant A_n$、$y_k\leqslant B_n$，所以
\[
  x_k+y_k\leqslant A_n+B_n.
\]
因此
\[
  \sup_{k\geqslant n}(x_k+y_k)
  \leqslant A_n+B_n.
\]
令 $n\to\infty$，利用公式 (3.2)，得到
\[
  \overline{\lim}_{n\to\infty}(x_n+y_n)
  \leqslant
  \overline{\lim}_{n\to\infty}x_n
  +\overline{\lim}_{n\to\infty}y_n.
\]
若某一上极限为 $+\infty$ 而和式仍有意义，不等式自动成立；其余允许的扩展实数情形按同一约定处理。
\end{xsolution}

\begin{xexercise}[3.6.4 练习题 11]
证明命题 3.6.5。
\end{xexercise}
\begin{xsolution}
记 $z_n=x_n+y_n$，并假设题中出现的和有意义。先由公式 (3.6) 应用于
\[
  x_n=z_n+(-y_n)
\]
得到
\[
  \underline{\lim}x_n
  \geqslant\underline{\lim}z_n+\underline{\lim}(-y_n)
  =\underline{\lim}z_n-\overline{\lim}y_n.
\]
移项得
\[
  \underline{\lim}(x_n+y_n)
  \leqslant\underline{\lim}x_n+\overline{\lim}y_n.
\]

再由公式 (3.5) 应用于
\[
  y_n=z_n+(-x_n)
\]
得到
\[
  \overline{\lim}y_n
  \leqslant\overline{\lim}z_n+\overline{\lim}(-x_n)
  =\overline{\lim}z_n-\underline{\lim}x_n.
\]
移项即
\[
  \underline{\lim}x_n+\overline{\lim}y_n
  \leqslant\overline{\lim}(x_n+y_n).
\]
合并两式正是命题 3.6.5。
\end{xsolution}

\section{对于教学的建议}
\subsection*{3.7.3 第一组参考题}

\begin{xreferenceproblem}[第一组参考题 1]
证明：数列有界的充分必要条件是它的每个子列有收敛子列。
\end{xreferenceproblem}
\begin{xsolution}
若 $\{x_n\}$ 有界，则它的任意子列仍然有界。对任意子列应用凝聚定理，都能从中再选出收敛子列。

反之，假设 $\{x_n\}$ 无界。递归选取下标
\[
  n_1<n_2<\cdots
\]
使
\[
  \abs{x_{n_k}}>k,\qquad k\in\NN_+.
\]
于是子列 $\{x_{n_k}\}$ 是无穷大量。它的任何子列仍有绝对值趋于 $+\infty$，因而不可能有收敛子列。这与题设“每个子列都有收敛子列”矛盾。因此原数列必有界。
\end{xsolution}

\begin{xreferenceproblem}[第一组参考题 2]
证明：数列收敛的充分必要条件是存在一个数 $a$，使数列的每个子列有收敛于 $a$ 的子列。
\end{xreferenceproblem}
\begin{xsolution}
若 $x_n\to a$，则任意子列也收敛于 $a$，当然它本身就是一个收敛于 $a$ 的子列。

反之，假设存在 $a$ 具有题设性质，但 $x_n$ 不收敛于 $a$。由极限定义的否定，存在 $\varepsilon_0>0$，对任意 $N$ 都有某个 $n>N$ 满足
\[
  \abs{x_n-a}\geqslant\varepsilon_0.
\]
由这些项可选出下标严格增加的子列 $\{x_{n_k}\}$，并且每一项都满足
\[
  \abs{x_{n_k}-a}\geqslant\varepsilon_0.
\]
这个子列的任何子列仍满足同样不等式，所以没有任何子列能够收敛于 $a$，与题设矛盾。故必有 $x_n\to a$。
\end{xsolution}

\begin{xreferenceproblem}[第一组参考题 3]
证明：在有界闭区间上的无界函数，一定在区间内某一点的每一个邻域上无界。又问：开区间上的无界函数是否有类似性质？
\end{xreferenceproblem}
\begin{xsolution}
设函数 $f$ 在闭区间 $[a,b]$ 上无界。反设对每个 $\xi\in[a,b]$，都存在邻域 $O(\xi)$，使 $f$ 在 $O(\xi)\cap[a,b]$ 上有界。于是对每个 $\xi$ 可取常数 $M_\xi$，使
\[
  \abs{f(x)}\leqslant M_\xi,
  \qquad x\in O(\xi)\cap[a,b].
\]
邻域族 $\{O(\xi)\}$ 构成 $[a,b]$ 的开覆盖。由覆盖定理，存在有限子覆盖
\[
  O(\xi_1),\ldots,O(\xi_m).
\]
令 $M=\max_{1\leqslant j\leqslant m}M_{\xi_j}$，则对每个 $x\in[a,b]$ 都有 $\abs{f(x)}\leqslant M$，这与 $f$ 无界矛盾。因此至少存在一点 $\xi\in[a,b]$，使 $f$ 在 $\xi$ 的每个邻域与 $[a,b]$ 的交上都无界。

对开区间结论一般不成立。例如在 $(0,1)$ 上取
\[
  f(x)=\frac1x.
\]
它在整个 $(0,1)$ 上无界，但对任意固定 $\xi\in(0,1)$，取 $0<\delta<\xi/2$，则在
\[
  (\xi-\delta,\xi+\delta)\subset(\xi/2,1)
\]
上有 $\abs{f(x)}<2/\xi$。所以它在每个内点的某个邻域上都有界，没有题述的“坏点”。
\end{xsolution}

\begin{xreferenceproblem}[第一组参考题 4]
设函数 $f$ 在 $(a,b)$ 上满足原题的局部严格增性质。证明 $f$ 在 $(a,b)$ 上严格单调增加。
\end{xreferenceproblem}
\begin{xsolution}
反设结论不成立，则存在 $u<v$，$u,v\in(a,b)$，使
\[
  f(u)\geqslant f(v).
\]
从区间 $I_1=[u,v]$ 开始构造闭区间套，并保持每个区间 $I_n=[u_n,v_n]$ 满足
\[
  f(u_n)\geqslant f(v_n).
\]
设中点为 $c_n=(u_n+v_n)/2$。若
\[
  f(u_n)\geqslant f(c_n),
\]
取 $I_{n+1}=[u_n,c_n]$；否则 $f(c_n)>f(u_n)\geqslant f(v_n)$，取 $I_{n+1}=[c_n,v_n]$。这样总能保持上述不等式，而且
\[
  \abs{I_n}=2^{1-n}(v-u)\to0.
\]
由闭区间套定理，存在唯一 $\xi\in[u,v]\subset(a,b)$，使
\[
  u_n\to\xi,\qquad v_n\to\xi.
\]
按题设，在 $\xi$ 处存在 $\delta>0$，使得在 $\xi$ 的这个邻域内，左侧点的函数值严格小于 $f(\xi)$，右侧点的函数值严格大于 $f(\xi)$。

取 $n$ 足够大，使 $u_n,v_n\in(\xi-\delta,\xi+\delta)$。若 $u_n<\xi<v_n$，则
\[
  f(u_n)<f(\xi)<f(v_n),
\]
与 $f(u_n)\geqslant f(v_n)$ 矛盾。若 $u_n=\xi<v_n$，则题设给出 $f(u_n)=f(\xi)<f(v_n)$，仍矛盾；若 $u_n<\xi=v_n$，则题设给出 $f(u_n)<f(\xi)=f(v_n)$，也与 $f(u_n)\geqslant f(v_n)$ 矛盾。因此不存在这样的 $u<v$，故对任意 $u<v$ 都有 $f(u)<f(v)$，即 $f$ 在 $(a,b)$ 上严格单调增加。
\end{xsolution}

\begin{xreferenceproblem}[第一组参考题 5]
试用上、下极限的工具证明 Stolz 定理。
\end{xreferenceproblem}
\begin{xsolution}
证明常用的 Stolz 形式：设 $\{y_n\}$ 严格增加且 $y_n\to+\infty$，令
\[
  d_n=\frac{x_{n+1}-x_n}{y_{n+1}-y_n}.
\]
则
\[
  \underline{\lim}_{n\to\infty}d_n
  \leqslant
  \underline{\lim}_{n\to\infty}\frac{x_n}{y_n}
  \leqslant
  \overline{\lim}_{n\to\infty}\frac{x_n}{y_n}
  \leqslant
  \overline{\lim}_{n\to\infty}d_n,
\]
只要相应的扩展实数运算有意义。因此特别地，若 $d_n\to L$，则 $x_n/y_n\to L$。

先证最右不等式。记
\[
  D=\overline{\lim}_{n\to\infty}d_n.
\]
先设 $D$ 有限。任取 $q>D$。由上极限的性质，存在 $N$，当 $k\geqslant N$ 时
\[
  d_k<q.
\]
由于 $y_{k+1}-y_k>0$，故
\[
  x_{k+1}-x_k<q(y_{k+1}-y_k).
\]
从 $k=N$ 到 $n-1$ 求和，得到
\[
  x_n-x_N<q(y_n-y_N),
\]
从而
\[
  \frac{x_n}{y_n}
  <q+\frac{x_N-qy_N}{y_n}.
\]
令 $n\to\infty$，因 $y_n\to+\infty$，有
\[
  \overline{\lim}_{n\to\infty}\frac{x_n}{y_n}\leqslant q.
\]
令 $q\downarrow D$，即得
\[
  \overline{\lim}\frac{x_n}{y_n}\leqslant\overline{\lim}d_n.
\]

再证最左不等式。记 $d=\underline{\lim}d_n$，先设 $d$ 有限。任取 $q<d$，存在 $N$，当 $k\geqslant N$ 时 $d_k>q$，即
\[
  x_{k+1}-x_k>q(y_{k+1}-y_k).
\]
从 $k=N$ 到 $n-1$ 求和，得到
\[
  x_n-x_N>q(y_n-y_N),
\]
故
\[
  \frac{x_n}{y_n}>q+\frac{x_N-qy_N}{y_n}.
\]
令 $n\to\infty$，得
\[
  \underline{\lim}_{n\to\infty}\frac{x_n}{y_n}\geqslant q.
\]
再令 $q\uparrow d$，即得
\[
  \underline{\lim}\frac{x_n}{y_n}\geqslant\underline{\lim}d_n.
\]
若相应上下极限为无穷，则对任意有限 $q$ 使用上述同一估计，便得到相应的扩展实数结论。
\end{xsolution}

\begin{xreferenceproblem}[第一组参考题 6]
设 $\{x_n\},\{y_n\}$ 为非负数列，证明原题所列关于乘积的四重上下极限不等式。
\end{xreferenceproblem}
\begin{xsolution}
先讨论所有有关上下极限均有限的情形。对每个 $n$ 记
\[
  \alpha_n=\inf_{k\geqslant n}x_k,
  \quad A_n=\sup_{k\geqslant n}x_k,
  \quad
  \beta_n=\inf_{k\geqslant n}y_k,
  \quad B_n=\sup_{k\geqslant n}y_k.
\]
由于数列非负，对每个 $k\geqslant n$ 有
\[
  \alpha_n\beta_n\leqslant x_ky_k\leqslant A_nB_n.
\]
故
\[
  \alpha_n\beta_n
  \leqslant\inf_{k\geqslant n}(x_ky_k),
  \qquad
  \sup_{k\geqslant n}(x_ky_k)\leqslant A_nB_n.
\]

还需证明中间两步。任取 $\varepsilon>0$。由下确界定义，可取 $k\geqslant n$ 使
\[
  x_k<\alpha_n+\varepsilon.
\]
同时 $y_k\leqslant B_n$，所以
\[
  \inf_{j\geqslant n}(x_jy_j)
  \leqslant x_ky_k
  \leqslant(\alpha_n+\varepsilon)B_n.
\]
令 $\varepsilon\downarrow0$，得
\[
  \inf_{k\geqslant n}(x_ky_k)\leqslant\alpha_nB_n.
\]
另一方面，可取 $k\geqslant n$ 使 $y_k>B_n-\varepsilon$；又 $x_k\geqslant\alpha_n$，故
\[
  \sup_{j\geqslant n}(x_jy_j)
  \geqslant x_ky_k
  \geqslant\alpha_n(B_n-\varepsilon).
\]
令 $\varepsilon\downarrow0$，得
\[
  \alpha_nB_n\leqslant\sup_{k\geqslant n}(x_ky_k).
\]

将这些不等式合并并令 $n\to\infty$，利用尾集表示，即得
\[
  \underline{\lim}x_n\,\underline{\lim}y_n
  \leqslant\underline{\lim}(x_ny_n)
  \leqslant\underline{\lim}x_n\,\overline{\lim}y_n
  \leqslant\overline{\lim}(x_ny_n)
  \leqslant\overline{\lim}x_n\,\overline{\lim}y_n.
\]
若出现 $+\infty$，只要题设所说的各乘积有意义，可用有限截断或任意大的有限界代替上述 $A_n,B_n$，再令截断值趋于无穷，得到同样结论。
\end{xsolution}

\begin{xreferenceproblem}[第一组参考题 7]
设 $\{x_n\}$ 为正数列。若 $\lim x_{n+1}/x_n=l$，证明 $\lim\sqrt[n]{x_n}=l$。
\end{xreferenceproblem}
\begin{xsolution}
由 3.6.4 练习题 5 已证明
\[
  \underline{\lim}_{n\to\infty}\frac{x_{n+1}}{x_n}
  \leqslant
  \underline{\lim}_{n\to\infty}\sqrt[n]{x_n}
  \leqslant
  \overline{\lim}_{n\to\infty}\sqrt[n]{x_n}
  \leqslant
  \overline{\lim}_{n\to\infty}\frac{x_{n+1}}{x_n}.
\]
题设比值数列收敛于 $l$，所以左右两端都等于 $l$。夹逼得到
\[
  \underline{\lim}_{n\to\infty}\sqrt[n]{x_n}
  =\overline{\lim}_{n\to\infty}\sqrt[n]{x_n}=l,
\]
因此
\[
  \lim_{n\to\infty}\sqrt[n]{x_n}=l.
\]
对于 $l=0$ 或 $l=+\infty$，上述扩展实数形式的夹逼仍给出相同结论。
\end{xsolution}

\begin{xreferenceproblem}[第一组参考题 8]
若数列 $\{a_n\}$ 的每个子列 $\{a_{n_k}\}$ 都满足
\[
  \lim_{k\to\infty}\frac{a_{n_1}+\cdots+a_{n_k}}{k}=a,
\]
证明 $a_n\to a$。
\end{xreferenceproblem}
\begin{xsolution}
反设 $a_n$ 不收敛于 $a$。则存在 $\varepsilon_0>0$，使集合
\[
  \{n\in\NN_+\mid \abs{a_n-a}\geqslant\varepsilon_0\}
\]
为无限集。在这些下标中，至少有一类是无限的：
\[
  a_n\geqslant a+\varepsilon_0
  \quad\text{或}\quad
  a_n\leqslant a-\varepsilon_0.
\]
若前者无限，选出相应子列 $\{a_{n_k}\}$，则对每个 $k$，
\[
  \frac{a_{n_1}+\cdots+a_{n_k}}k\geqslant a+\varepsilon_0,
\]
其极限不可能是 $a$。若后者无限，选出所有满足 $a_{n_k}\leqslant a-\varepsilon_0$ 的相应子列，则对每个 $k$ 都有 $k$ 项平均不超过 $a-\varepsilon_0$，其极限也不可能是 $a$。两种情况都与题设矛盾，故 $a_n\to a$。
\end{xsolution}

\begin{xreferenceproblem}[第一组参考题 9]
设 $\{x_n\}$ 为正数列，证明原扫描题中的不等式，并说明常数 $1$ 为最佳值。
\end{xreferenceproblem}
% SOURCE-NOTE: 原扫描页（上册物理页 112，印刷页 95）本题为上极限；chapters/ch03.tex 误录为下极限。此处按原扫描题意作答，未修改章节正文。
\begin{xsolution}
按原扫描题意，需证明
\[
  \overline{\lim}_{n\to\infty}
  n\left(\frac{1+x_{n+1}}{x_n}-1\right)\geqslant1.
\]
记
\[
  u_n=n\left(\frac{1+x_{n+1}}{x_n}-1\right).
\]
反设 $\overline{\lim}u_n<1$。则可取 $\delta\in(0,1)$ 和 $N$，使当 $n\geqslant N$ 时
\[
  u_n\leqslant1-\delta.
\]
这等价于
\[
  x_{n+1}\leqslant
  \left(1+\frac{1-\delta}{n}\right)x_n-1.
\]
令 $v_n=x_n/n>0$，两边除以 $n+1$ 得
\[
  v_{n+1}
  \leqslant
  \left(1-\frac{\delta}{n+1}\right)v_n-\frac1{n+1}.
\]
再令
\[
  w_n=v_n+\frac1\delta.
\]
则
\[
  w_{n+1}\leqslant
  \left(1-\frac{\delta}{n+1}\right)w_n.
\]
于是对 $n>N$，
\[
  0<w_n
  \leqslant w_N\prod_{k=N+1}^{n}\left(1-\frac\delta k\right)
  \leqslant w_N\exp\left(-\delta\sum_{k=N+1}^{n}\frac1k\right)
  \longrightarrow0.
\]
但由 $v_n>0$ 有 $w_n>1/\delta$，矛盾。因此
\[
  \overline{\lim}u_n\geqslant1.
\]

下面证明 $1$ 不能增大。取
\[
  x_n=n\ln(n+1)>0.
\]
则
\begin{align*}
  u_n
  &=\frac{n}{x_n}+n\left(\frac{x_{n+1}}{x_n}-1\right)\\
  &=\frac1{\ln(n+1)}
    +1+
    \frac{(n+1)\ln\left(1+\frac1{n+1}\right)}{\ln(n+1)}.
\end{align*}
其中两个分式都趋于 $0$，故 $u_n\to1$。所以普遍下界不可能取比 $1$ 更大的常数，$1$ 为最佳值。
\end{xsolution}

\begin{xreferenceproblem}[第一组参考题 10]
设 $\{x_n\}$ 为正数列，证明
\[
  \overline{\lim}_{n\to\infty}
  \left(\frac{x_1+x_{n+1}}{x_n}\right)^n\geqslant\ee,
\]
并证明 $\ee$ 为最佳值。
\end{xreferenceproblem}
\begin{xsolution}
令
\[
  z_n=\frac{x_n}{x_1}>0,
\]
则 $z_1=1$，且
\[
  \frac{x_1+x_{n+1}}{x_n}
  =\frac{1+z_{n+1}}{z_n}.
\]
记
\[
  u_n=n\left(\frac{1+z_{n+1}}{z_n}-1\right).
\]
由上一题，$\overline{\lim}u_n\geqslant1$。任取 $\varepsilon\in(0,1)$，有无穷多个 $n$ 满足
\[
  u_n>1-\varepsilon.
\]
沿这些下标，
\[
  \left(\frac{1+z_{n+1}}{z_n}\right)^n
  =\left(1+\frac{u_n}{n}\right)^n
  \geqslant
  \left(1+\frac{1-\varepsilon}{n}\right)^n.
\]
令这些下标趋于无穷，右端趋于 $\ee^{1-\varepsilon}$，故
\[
  \overline{\lim}_{n\to\infty}
  \left(\frac{x_1+x_{n+1}}{x_n}\right)^n
  \geqslant\ee^{1-\varepsilon}.
\]
再令 $\varepsilon\downarrow0$，得到所需下界 $\ee$。

为说明最佳性，取
\[
  z_n=\frac{n\ln(n+1)}{\ln2},
\]
则 $z_1=1$，并由上一题的计算知
\[
  n\left(\frac{1+z_{n+1}}{z_n}-1\right)\to1.
\]
因此若记括号内为 $1+u_n/n$，则 $u_n\to1$，从而
\[
  \left(\frac{1+z_{n+1}}{z_n}\right)^n
  =\left(1+\frac{u_n}{n}\right)^n\to\ee.
\]
令 $x_n=x_1z_n$ 即得到正数列的例子，其上极限恰为 $\ee$。故 $\ee$ 为最佳值。
\end{xsolution}

\subsection*{3.7.3 第二组参考题}

\begin{xreferenceproblem}[第二组参考题 1]
从确界存在定理证明 Archimedes 原理：对任意正数 $a,b$，存在正整数 $n$ 使 $na>b$。
\end{xreferenceproblem}
\begin{xsolution}
反设对所有 $n\in\NN_+$ 都有 $na\leqslant b$。于是数集
\[
  S=\{na\mid n\in\NN_+\}
\]
非空且有上界。由确界存在定理，存在
\[
  \alpha=\sup S.
\]
由于 $a>0$，有 $\alpha-a<\alpha$，所以 $\alpha-a$ 不可能仍是 $S$ 的上界。因此存在某个 $n$ 使
\[
  na>\alpha-a.
\]
两边加 $a$ 得
\[
  (n+1)a>\alpha.
\]
但 $(n+1)a\in S$，这与 $\alpha$ 是 $S$ 的上界矛盾。故反设不成立，必存在正整数 $n$ 使 $na>b$。
\end{xsolution}

\begin{xreferenceproblem}[第二组参考题 2]
证明 Dedekind 连续性结论：若 $\RR=A\cup B$，且每个 $A$ 中的数都小于每个 $B$ 中的数，则或者 $A$ 有最大数而 $B$ 无最小数，或者 $B$ 有最小数而 $A$ 无最大数。
\end{xreferenceproblem}
\begin{xsolution}
任取 $b_0\in B$。由题设，每个 $a\in A$ 都满足 $a<b_0$，所以 $A$ 有上界。令
\[
  c=\sup A.
\]
因为 $\RR=A\cup B$，必有 $c\in A$ 或 $c\in B$。此外 $A\cap B=\varnothing$，否则若 $c\in A\cap B$，题设将给出 $c<c$。

若 $c\in A$，则由 $c$ 是 $A$ 的上界可知 $c=\max A$。此时 $B$ 不可能有最小数。否则设 $d=\min B$，由题设 $c<d$。取中点
\[
  m=\frac{c+d}{2},
\]
则 $c<m<d$。若 $m\in A$，与 $c=\max A$ 矛盾；若 $m\in B$，与 $d=\min B$ 矛盾。故 $B$ 无最小数。

若 $c\in B$，则对每个 $b\in B$，因为 $b$ 是 $A$ 的上界，有 $c=\sup A\leqslant b$，所以 $c=\min B$。若 $A$ 有最大数 $d$，则 $d<c$；中点 $(d+c)/2$ 属于 $A$ 或 $B$，分别与 $d$ 最大或 $c$ 最小矛盾。因此 $A$ 无最大数。

故恰有题述两种情形之一成立。
\end{xsolution}

\begin{xreferenceproblem}[第二组参考题 3]
证明实数的连通性：把 $\RR$ 分成两个非空集合 $A,B$，则或者 $A$ 中有数列收敛于 $B$ 中的点，或者 $B$ 中有数列收敛于 $A$ 中的点。
\end{xreferenceproblem}
\begin{xsolution}
任取 $a_1\in A$ 与 $b_1\in B$。以这两点为端点得到闭区间 $I_1$。设已经得到闭区间 $I_n$，它的两个端点分别属于 $A$ 与 $B$。取中点 $c_n$。由于 $A\cup B=\RR$，$c_n$ 属于 $A$ 或 $B$。保留 $c_n$ 与原来属于另一个集合的端点，得到一个半区间 $I_{n+1}$。于是 $I_{n+1}\subset I_n$，它的两个端点仍分别属于 $A,B$，且
\[
  \abs{I_{n+1}}=\frac12\abs{I_n}.
\]

由闭区间套定理，存在唯一
\[
  \xi\in\bigcap_{n=1}^{\infty}I_n.
\]
对每个 $n$，把 $I_n$ 中属于 $A$ 的端点记为 $\alpha_n$，属于 $B$ 的端点记为 $\beta_n$。因为 $\xi\in I_n$，
\[
  \abs{\alpha_n-\xi}\leqslant\abs{I_n}\to0,
  \qquad
  \abs{\beta_n-\xi}\leqslant\abs{I_n}\to0.
\]
故
\[
  \alpha_n\to\xi,
  \qquad
  \beta_n\to\xi.
\]
若 $\xi\in A$，则 $B$ 中的数列 $\{\beta_n\}$ 收敛于 $A$ 中的点 $\xi$；若 $\xi\in B$，则 $A$ 中的数列 $\{\alpha_n\}$ 收敛于 $B$ 中的点 $\xi$。结论得证。
\end{xsolution}

\begin{xreferenceproblem}[第二组参考题 4]
用压缩映射原理证明原题所列交错嵌套根式数列收敛，并求其极限。
\end{xreferenceproblem}
\begin{xsolution}
把原数列记为 $\{x_n\}$。从根式结构可直接看出，每隔两项满足递推关系
\[
  x_{n+2}=F(x_n),
  \qquad
  F(x)=\sqrt{7-\sqrt{7+x}}.
\]
前两项为
\[
  x_1=\sqrt7,
  \qquad
  x_2=\sqrt{7-\sqrt7},
\]
都属于 $[1,3]$。当 $x\in[1,3]$ 时，
\[
  F(x)\in
  \left[\sqrt{7-\sqrt{10}},\sqrt{7-\sqrt8}\right]
  \subset[1,3],
\]
所以 $F$ 把 $[1,3]$ 映入自身。

对 $x,y\in[1,3]$，两次有理化得到
\begin{align*}
  \abs{F(x)-F(y)}
  &=\frac{\abs{\sqrt{7+x}-\sqrt{7+y}}}{F(x)+F(y)}\\
  &=\frac{\abs{x-y}}
  {(\sqrt{7+x}+\sqrt{7+y})(F(x)+F(y))}.
\end{align*}
在 $[1,3]$ 上有
\[
  \sqrt{7+x}+\sqrt{7+y}\geqslant2\sqrt8=4\sqrt2,
\]
以及
\[
  F(x)+F(y)\geqslant2\sqrt{7-\sqrt{10}}.
\]
两者乘积为
\[
  8\sqrt{2(7-\sqrt{10})}>8,
\]
故可取一个简单的统一估计
\[
  \abs{F(x)-F(y)}\leqslant\frac18\abs{x-y}.
\]
因此 $F$ 是 $[1,3]$ 上的压缩映射。

于是奇数子列
\[
  x_1,x_3,x_5,\ldots
\]
和偶数子列
\[
  x_2,x_4,x_6,\ldots
\]
都是同一压缩映射 $F$ 的迭代序列，故二者都收敛到 $F$ 在 $[1,3]$ 中唯一的不动点。直接计算
\[
  F(2)=\sqrt{7-\sqrt9}=2,
\]
所以该唯一不动点为 $2$。因奇、偶两个子列都趋于 $2$，整个数列也收敛，而且
\[
  \boxed{\lim_{n\to\infty}x_n=2}.
\]
\end{xsolution}

\begin{xreferenceproblem}[第二组参考题 5]
设有界数列 $\{x_n\}$ 对每个数列 $\{y_n\}$ 都满足
\[
  \overline{\lim}(x_n+y_n)
  =\overline{\lim}x_n+\overline{\lim}y_n.
\]
证明 $\{x_n\}$ 收敛。
\end{xreferenceproblem}
\begin{xsolution}
由于 $\{x_n\}$ 有界，记
\[
  L=\overline{\lim}_{n\to\infty}x_n,
  \qquad
  l=\underline{\lim}_{n\to\infty}x_n,
\]
二者都是有限数。题设对每个 $\{y_n\}$ 成立，特别取
\[
  y_n=-x_n.
\]
则左边恒为
\[
  \overline{\lim}_{n\to\infty}(x_n+y_n)=0.
\]
另一方面
\[
  \overline{\lim}_{n\to\infty}y_n
  =\overline{\lim}_{n\to\infty}(-x_n)=-l.
\]
所以题设恒等式给出
\[
  0=L-l.
\]
因此 $L=l$。有界数列的上、下极限相等当且仅当它收敛，所以 $\{x_n\}$ 收敛。
\end{xsolution}

\begin{xreferenceproblem}[第二组参考题 6]
证明原题的两个关于正数列的结论。
\end{xreferenceproblem}
\begin{xsolution}
\begin{enumerate}[label=(\arabic*)]
  \item 设 $x_n>0$ 且 $x_n\to0$。称 $n$ 为一个“新低点”下标，如果
  \[
    x_n<x_k,\qquad k=1,2,\ldots,n-1.
  \]
  反设这样的下标只有有限多个，令 $n_0$ 为最后一个新低点下标。则 $x_{n_0}$ 是前 $n_0$ 项中的最小值。以后也不可能出现 $x_n<x_{n_0}$，否则第一次出现这种情形的下标就会成为新的新低点，与 $n_0$ 最后矛盾。因此对所有 $n>n_0$ 都有
  \[
    x_n\geqslant x_{n_0}>0,
  \]
  这与 $x_n\to0$ 矛盾。故这样的 $n$ 有无限多个。

  \item 设 $x_n>0$ 且存在 $c>0$ 使 $x_n\geqslant c$。反设
  \[
    R=\overline{\lim}_{n\to\infty}\frac{x_{n+1}}{x_n}<1.
  \]
  可取 $q$ 满足 $R<q<1$。由上极限定义，存在 $N$，当 $n\geqslant N$ 时
  \[
    \frac{x_{n+1}}{x_n}<q.
  \]
  归纳得到
  \[
    x_{N+k}\leqslant x_Nq^k\to0,
  \]
  这与 $x_{N+k}\geqslant c>0$ 矛盾。因此
  \[
    \overline{\lim}_{n\to\infty}\frac{x_{n+1}}{x_n}\geqslant1.
  \]
\end{enumerate}
\end{xsolution}

\begin{xreferenceproblem}[第二组参考题 7]
设 $y_n=px_n+qx_{n+1}$，其中 $\abs p<\abs q$。证明：若 $\{y_n\}$ 收敛，则 $\{x_n\}$ 也收敛。
\end{xreferenceproblem}
\begin{xsolution}
设
\[
  y_n\to C.
\]
由于 $\abs p<\abs q$，有 $q\neq0$ 且 $p+q\neq0$。令
\[
  L=\frac{C}{p+q},
  \qquad z_n=x_n-L.
\]
由
\[
  qx_{n+1}=y_n-px_n
\]
以及 $(p+q)L=C$，可得
\[
  z_{n+1}=r z_n+\delta_n,
  \qquad
  r=-\frac pq,
  \qquad
  \delta_n=\frac{y_n-C}{q}.
\]
这里
\[
  \abs r<1,
  \qquad \delta_n\to0.
\]

任取 $\varepsilon>0$。取 $N$ 使当 $n\geqslant N$ 时
\[
  \abs{\delta_n}<\frac{1-\abs r}{2}\varepsilon.
\]
递推估计给出对 $k\geqslant1$，
\begin{align*}
  \abs{z_{N+k}}
  &\leqslant \abs r^k\abs{z_N}
  +\sum_{j=0}^{k-1}\abs r^j
  \frac{1-\abs r}{2}\varepsilon\\
  &\leqslant \abs r^k\abs{z_N}+\frac\varepsilon2.
\end{align*}
取 $k$ 足够大，使第一项小于 $\varepsilon/2$，便有 $\abs{z_{N+k}}<\varepsilon$。故 $z_n\to0$，即
\[
  x_n\to L=\frac{C}{p+q}.
\]
所以 $\{x_n\}$ 收敛。
\end{xsolution}

\begin{xreferenceproblem}[第二组参考题 8]
设 $\{x_n\}$ 有界，且 $x_{2n}+2x_n\to A$。证明 $\{x_n\}$ 收敛，并求极限。
\end{xreferenceproblem}
\begin{xsolution}
记
\[
  L=\overline{\lim}_{n\to\infty}x_n,
  \qquad
  l=\underline{\lim}_{n\to\infty}x_n.
\]
因原数列有界，$L,l$ 均有限。令
\[
  y_n=x_{2n}+2x_n\to A.
\]
由
\[
  x_{2n}=y_n-2x_n
\]
及收敛数列与上下极限的运算法则，得到
\[
  \overline{\lim}_{n\to\infty}x_{2n}=A-2l,
  \qquad
  \underline{\lim}_{n\to\infty}x_{2n}=A-2L.
\]
另一方面，$\{x_{2n}\}$ 是原数列的子列，所以
\[
  \overline{\lim}x_{2n}\leqslant L,
  \qquad
  \underline{\lim}x_{2n}\geqslant l.
\]
故
\[
  A-2l\leqslant L,
  \qquad
  A-2L\geqslant l.
\]
即
\[
  A\leqslant L+2l,
  \qquad
  A\geqslant2L+l.
\]
于是
\[
  2L+l\leqslant L+2l,
\]
从而 $L\leqslant l$。而总有 $l\leqslant L$，所以 $L=l$，故 $\{x_n\}$ 收敛。设极限为 $s$，在题设关系中令 $n\to\infty$，得到
\[
  s+2s=A,
\]
所以
\[
  \boxed{s=\frac A3}.
\]
\end{xsolution}

\begin{xreferenceproblem}[第二组参考题 9]
设 $x_n=\sin n$。证明数列 $\{x_n\}$ 的极限点集合为 $[-1,1]$。
\end{xreferenceproblem}
\begin{xsolution}
先用一个标准的无理旋转事实。记
\[
  \langle t\rangle=t-\lfloor t\rfloor
\]
为 $t$ 的小数部分。若 $\alpha$ 为无理数，则集合
\[
  \{\langle n\alpha\rangle\mid n\in\NN_+\}
\]
在 $[0,1]$ 中稠密，而且每个开区间中有无穷多个这样的点。

给出证明。任取正整数 $m$。考察 $m+1$ 个小数部分
\[
  0,\langle\alpha\rangle,\langle2\alpha\rangle,\ldots,\langle m\alpha\rangle.
\]
把 $[0,1]$ 分成 $m$ 个等长小区间，由抽屉原理，有两个小数部分的距离小于 $1/m$。因此存在正整数 $q\leqslant m$，使 $q\alpha$ 到某个整数的距离
\[
  0<\delta<\frac1m.
\]
于是 $q\alpha$ 的小数部分或者等于 $\delta$，或者等于 $1-\delta$。相应的正整数倍 $jq\alpha$ 的小数部分按步长 $\delta$ 正向或反向绕单位圆排列，因此离任意给定的 $t\in[0,1]$ 都可小于 $1/m$。令 $m\to\infty$，便得稠密性。又从任意下标 $N$ 后开始的集合
\[
  \{\langle(N+n)\alpha\rangle\mid n\in\NN_+\}
\]
只是上述稠密集合在单位圆上的一个平移，仍然稠密，故每个开区间实际上被命中无穷多次。

现在取
\[
  \alpha=\frac1{2\pi}.
\]
由于 $\pi$ 为无理数，$\alpha$ 也为无理数，所以 $n$ 模 $2\pi$ 的余数在 $[0,2\pi]$ 中稠密并且无限次进入任意开子区间。

任取 $y\in[-1,1]$，选取 $\theta\in[0,2\pi]$ 使
\[
  \sin\theta=y.
\]
由上述稠密性，可选严格增加的正整数 $n_k$，使 $n_k$ 模 $2\pi$ 的余数趋于 $\theta$。由正弦函数连续性，
\[
  \sin n_k\to\sin\theta=y.
\]
因此 $[-1,1]$ 中每个点都是 $\{\sin n\}$ 的极限点。

反过来，因为对所有 $n$ 都有 $-1\leqslant\sin n\leqslant1$，任何收敛子列的极限也只能属于闭区间 $[-1,1]$。故极限点集合恰为
\[
  \boxed{[-1,1]}.
\]
\end{xsolution}

\begin{xreferenceproblem}[第二组参考题 10]
设 $\{x_n\}$ 有界，且 $x_{n+1}-x_n\to0$。若下极限和上极限分别为 $l,L$，证明 $[l,L]$ 中每一点都是极限点。
\end{xreferenceproblem}
\begin{xsolution}
由于数列有界，$l,L$ 都是有限极限点，所以端点 $l,L$ 本身已经是极限点。若 $l=L$，结论立即成立。以下设 $l<L$，任取
\[
  c\in(l,L).
\]

由 $l<c<L$ 和上下极限的定义，数列中满足 $x_n<c$ 的项有无穷多个，满足 $x_n>c$ 的项也有无穷多个，而且两类项都在任意大的下标以后出现。因此可以找到下标趋于无穷的一系列“跨越” $c$ 的相邻项，即存在 $n_k\to\infty$，使每个 $k$ 都有
\[
  x_{n_k}\leqslant c<x_{n_k+1}
  \quad\text{或}\quad
  x_{n_k}\geqslant c>x_{n_k+1}.
\]
具体地，从任意一个位于 $c$ 下方的充分晚项出发，向后寻找下一项位于 $c$ 上方的项，在第一次越过 $c$ 时便出现这样的相邻对；若从位于 $c$ 上方的充分晚项出发并向后寻找下一项位于 $c$ 下方的项，在第一次下降越过 $c$ 时也得到这样的相邻对。由于上下两类项都无限次出现，这一过程可反复进行并使 $n_k$ 严格增加。

在每一对跨越项中，选取离 $c$ 较近的一项，记为 $x_{m_k}$。则
\[
  \abs{x_{m_k}-c}
  \leqslant\abs{x_{n_k+1}-x_{n_k}}.
\]
又因 $n_k\to\infty$ 且
\[
  x_{n+1}-x_n\to0,
\]
右端趋于 $0$，所以
\[
  x_{m_k}\to c.
\]
故任意 $c\in(l,L)$ 都是数列的极限点。连同端点 $l,L$，得到 $[l,L]$ 中每一个点都是 $\{x_n\}$ 的极限点。
\end{xsolution}
